% GENERATED FILE. Edit canonical lessons, then run npm run generate:books.
% canonical-source-hash: fnv1a64:ab42532b4724f805
% canonical-lessons: FA-C294-bebakhshid, FA-C294-dast-e-rast, FA-R294-asking-the-way

\chapter{Asking The Way}
\label{ch:fa-asking-the-way}
\hlchaptermodality{294}

\begin{chapteropening}
\textbf{By the end of this chapter:} \emph{I can stop someone politely, ask how to get somewhere, and follow directions to the right, to the left and straight on.}

Complete FA-R294-asking-the-way: ask the way to the station and repeat the directions back.
\end{chapteropening}

\section[\texorpdfstring{\fa{ببخشید} (bebakhshid)}{ببخشید (bebakhshid)}]{\fa{ببخشید} (bebakhshid) — excuse me, and the way to the station}

\label{lesson:FA-C294-bebakhshid}



\begin{quote}
Say \textquotedblleft{}to forgive\textquotedblright{}, then \textquotedblleft{}where?\textquotedblright{}. (\emph{Bakhshidan}, \emph{kojâ}.) Then \textquotedblleft{}I want to go\textquotedblright{}. (\emph{Mikhâham beravam.})

On the street the first word you need is the one that stops a stranger politely.
\end{quote}

\subsection*{You'll want to know: \fa{ببخشید}}
\textbf{\fa{ببخشید}} — \emph{bebakhshid} — \textquotedblleft{}excuse me\textquotedblright{}, before a question; \textquotedblleft{}sorry\textquotedblright{}, after a bump.

It is \emph{bakhshidan}, \textquotedblleft{}to forgive\textquotedblright{}, with \emph{be-} and the \emph{shomâ} ending, the shape a polite request takes: \textquotedblleft{}forgive\textquotedblright{}. Persian stops a stranger by asking pardon first.

\begin{quote}
\textbf{\fa{ببخشید،} \fa{ایستگاه} \fa{کجاست؟}} — \emph{bebakhshid, istgâh kojâst?} — Excuse me, where is the station?
\end{quote}

\emph{Kojâst} is \emph{kojâ} with \emph{ast} fused on, the way \emph{chi} and \emph{ast} became \emph{chist}. The answer may start \textbf{\fa{نزدیک} \fa{است}}, \emph{nazdik ast}, \textquotedblleft{}it's near\textquotedblright{}.

\begin{grammarlens}[title={Grammar Lens: chetor mitavânam be ... beravam?}]
\emph{Tavânestan}, \textquotedblleft{}to be able\textquotedblright{}, has the present stem \textbf{tavân-}: \emph{mitavânam}, I can; \emph{mitavânid}, you can. Like \emph{mikhâham}, it takes the subjunctive after it. Put \emph{chetor}, \textquotedblleft{}how\textquotedblright{}, in front and you have the way:

\begin{quote}
\textbf{\fa{ببخشید،} \fa{چطور} \fa{میتوانم} \fa{به} \fa{ایستگاه} \fa{بروم؟}} — \emph{bebakhshid, chetor mitavânam be istgâh beravam?} — Excuse me, how can I get to the station?
\end{quote}

Swap the place and the question still works: \emph{be muze}, to the museum; \emph{be hotel}, to the hotel.
\end{grammarlens}

\begin{cousinweb}
\emph{bakhshidan} is native Persian. EtymDB 2.1 takes it back to a Middle Persian verb written \emph{bhš-} in Pahlavi letters, and lists it in two senses, \textquotedblleft{}to give, to grant\textquotedblright{} and \textquotedblleft{}to forgive\textquotedblright{}. Forgiving is a kind of granting, and \emph{bebakhshid} asks for it in advance.
\end{cousinweb}

\subsection*{Your turn}
Say these aloud:

\begin{itemize}
\raggedright
  \item \emph{bebakhshid, istgâh kojâst?}
  \item \emph{bebakhshid, chetor mitavânam be istgâh beravam?}
  \item the same question for the museum — \emph{be muze}
  \item the thanks after the answer — \emph{mamnun}
\end{itemize}

\begin{tcolorbox}[breakable,colback=teal!4,colframe=teal!35!black,title={Before you move on}]
What does \emph{bebakhshid} say, word for word? (\textbf{Forgive}, said politely.) \textquotedblleft{}Where is the station?\textquotedblright{} (\emph{\textbf{Istgâh kojâst?}}) \textquotedblleft{}How can I get to the station?\textquotedblright{} (\emph{\textbf{Chetor mitavânam be istgâh beravam?}})
\end{tcolorbox}
\section[\texorpdfstring{\fa{دست} \fa{راست} (dast-e râst)}{دست راست (dast-e râst)}]{\fa{دست} \fa{راست} (dast-e râst) — on the right, and how directions are given}

\label{lesson:FA-C294-dast-e-rast}



\begin{quote}
Say \textquotedblleft{}right\textquotedblright{}, \textquotedblleft{}left\textquotedblright{} and \textquotedblleft{}straight ahead\textquotedblright{}. (\emph{Râst}, \emph{chap}, \emph{mostaqim}.) Then ask the way to the station. (\emph{Bebakhshid, chetor mitavânam be istgâh beravam?})

Now the answer, which arrives fast, so it pays to know its shape.
\end{quote}

\subsection*{You'll want to know: \fa{دست} \fa{راست}}
\textbf{\fa{دست} \fa{راست}} — \emph{dast-e râst} — \textquotedblleft{}on the right\textquotedblright{}: word for word, \textquotedblleft{}the right hand\textquotedblright{}. \textbf{\fa{دست} \fa{چپ}}, \emph{dast-e chap}, is \textquotedblleft{}on the left\textquotedblright{}.

\begin{quote}
\textbf{\fa{ایستگاه} \fa{دست} \fa{راست} \fa{است}.} — \emph{istgâh dast-e râst ast} — The station is on the right.
\end{quote}

The ezafe joins the hand to its side, the same \emph{-e} as in \emph{esm-e man}.

\begin{grammarlens}[title={Grammar Lens: the polite instruction}]
A stranger gives directions in the subjunctive with the \emph{shomâ} ending, said as an instruction. You hold the shape already:

\begin{quote}
\textbf{\fa{مستقیم} \fa{بروید}.} — \emph{mostaqim beravid} — Go straight on.
\end{quote}

\begin{quote}
\textbf{\fa{سر} \fa{چهارراه} \fa{به} \fa{دست} \fa{چپ} \fa{بپیچید}.} — \emph{sar-e chahârrâh be dast-e chap bepichid} — At the crossroads, turn left.
\end{quote}

\emph{Bepichid} is \emph{pichidan}, \textquotedblleft{}to turn\textquotedblright{}, from chapter one hundred and thirty-three, on its stem \emph{pich-}. \emph{Sar-e chahârrâh} is \textquotedblleft{}at the head of the crossroads\textquotedblright{}. To a friend, someone you call \emph{to}, the ending drops: \emph{bepich}, turn.

So a whole answer runs: \emph{mostaqim beravid, ba'd be dast-e râst bepichid. istgâh dast-e chap ast.}
\end{grammarlens}

\begin{cousinweb}
Both words are inherited. \emph{dast}, \textquotedblleft{}hand\textquotedblright{}, continues Old Persian \emph{dasta}, on the same Indo-European root as Sanskrit \emph{hasta}, \textquotedblleft{}hand\textquotedblright{} (EtymDB 2.1). \emph{râst} continues Old Persian \emph{rāsta}, \textquotedblleft{}right, straight\textquotedblright{}, and EtymDB sets it beside English \textbf{right}: both go back to an Indo-European word for \textquotedblleft{}moved in a straight line\textquotedblright{}.
\end{cousinweb}

\subsection*{Your turn}
Say these aloud:

\begin{itemize}
\raggedright
  \item \emph{istgâh dast-e râst ast}, then the same on the left
  \item \textquotedblleft{}Go straight on\textquotedblright{} — \emph{mostaqim beravid}
  \item \emph{sar-e chahârrâh be dast-e chap bepichid}
  \item the whole answer back, as you would to check it: \emph{mostaqim, ba'd dast-e râst...}
\end{itemize}

\begin{tcolorbox}[breakable,colback=teal!4,colframe=teal!35!black,title={Before you move on}]
\textquotedblleft{}On the left\textquotedblright{}? (\emph{\textbf{Dast-e chap.}}) \textquotedblleft{}Turn right\textquotedblright{}, politely? (\emph{\textbf{Be dast-e râst bepichid.}}) Which English word does EtymDB set beside \emph{râst}? (\textbf{Right.})
\end{tcolorbox}
\section[(review)]{Review — on the street, asking the way}

\label{lesson:FA-R294-asking-the-way}



\begin{quote}
Stop someone and ask the way to the museum. (\emph{Bebakhshid, chetor mitavânam be muze beravam?})
\end{quote}

\begin{grammarlens}[title={What you've built}]
Two lessons, and a whole exchange. You open with \emph{bebakhshid}, ask with \emph{chetor mitavânam be ... beravam?} or \emph{... kojâst?}, and the answer comes back in the polite instruction: \emph{mostaqim beravid}, \emph{be dast-e chap bepichid}.

Here is one exchange, start to finish:

\begin{quote}
— \emph{bebakhshid, muze kojâst?} — \emph{mostaqim beravid. sar-e chahârrâh be dast-e râst bepichid. muze dast-e chap ast.} — \emph{mostaqim, ba'd dast-e râst, ba'd dast-e chap. mamnun.}
\end{quote}

The third line is the habit worth keeping: say the directions back, and any mistake shows up while the stranger is still there.
\end{grammarlens}

\subsection*{Your turn}
Say these aloud:

\begin{itemize}
\raggedright
  \item the three lines of the exchange, playing both people
  \item the same, for the station
  \item a reason you are lost — \emph{dalil in ast ke naqshe nadâram}, the reason is that I don't have a map
  \item advice to a visitor — \emph{bâyad mostaqim beravid}
\end{itemize}

\begin{tcolorbox}[breakable,colback=teal!4,colframe=teal!35!black,title={Before you move on}]
What do you say before the question? (\emph{\textbf{Bebakhshid.}}) \textquotedblleft{}Turn left\textquotedblright{}? (\emph{\textbf{Be dast-e chap bepichid.}}) Why say the directions back? (So a mistake \textbf{shows up at once}.)
\end{tcolorbox}
